\documentclass[french, a4paper, 12pt, titlepage]{article}
%% Peut remplacer "article" par "scrartcl" %%

\usepackage{a4wide}
%\usepackage[top=2cm, bottom=2cm, left=2cm, right=2cm]{geometry}
\raggedbottom % prevents vertical white space on pages that cannot be filled properly

\usepackage{hyperref}
\hypersetup{
	colorlinks=true,       	% false: boxed links; true: colored links
	linkcolor=black,          	% color of internal links
	urlcolor=blue,           	% color of external links
	citecolor=grey
}

\usepackage[T1]{fontenc}
%\usepackage{fourier}
%\usepackage{utopia}
%\usepackage{palatino}

\usepackage{lmodern}
%% ajouter fonte petite capitale grasse à lmodern avec celle de computer modern %%
\rmfamily
\DeclareFontShape{T1}{lmr}{b}{sc}{<->ssub*cmr/bx/sc}{}
\DeclareFontShape{T1}{lmr}{bx}{sc}{<->ssub*cmr/bx/sc}{}
%% /ajout %%
\usepackage{wrapfig}

%\usepackage[a4paper]{geometry} % marges plus petites que a4paper standard
\usepackage{listings} % insérer code source
\lstloadlanguages{sh,bash,awk,make}
%\usepackage{algorithm} % algorithmique
%\usepackage{algorithmic}
\usepackage{url}
\usepackage[usenames, dvipsnames]{color} % couleurs (nombre de base étendu)
\usepackage{graphicx} % insérer images
\usepackage[utf8]{inputenc}
\usepackage[french]{babel}
\usepackage{amsmath}
\usepackage{amsfonts}
\usepackage{amssymb}
\usepackage{amsthm}
\usepackage{multicol}
\definecolor{grey}{rgb}{0.96,0.96,0.96}
\definecolor{grey2}{rgb}{0.3,0.3,0.3}

%% Define listings params %%
\lstset{
%	numbers=left,
	language=bash,
	tabsize=4,
	frame=single, % cadre autour du code
	breaklines=true, % autorise couper ligne trop longue
	basicstyle=\small\ttfamily,
	numberstyle=\scriptsize\ttfamily,
	backgroundcolor=\color{grey},
	showstringspaces=false,
	keywordstyle=\color{OliveGreen},
	stringstyle=\color{BrickRed},
	commentstyle=\color{grey2}\it,
	stepnumber=1 % numérote toute les x lignes
}
% listing utf8 fr %
\lstset{%
	inputencoding=utf8,
	extendedchars=true,
	literate=
		{é}{{\'{e}}}1
		{è}{{\`{e}}}1
		{ê}{{\^{e}}}1
		{ë}{{\¨{e}}}1
		{û}{{\^{u}}}1
		{ù}{{\`{u}}}1
		{â}{{\^{a}}}1
		{à}{{\`{a}}}1
		{î}{{\^{i}}}1
		{ç}{{\c{c}}}1
		{Ç}{{\c{C}}}1
		{É}{{\'{E}}}1
		{Ê}{{\^{E}}}1
		{À}{{\`{A}}}1
		{Â}{{\^{A}}}1
		{Î}{{\^{I}}}1
}
%% /Define listings params %%

%% Francisation des algorithmes
%\renewcommand{\algorithmicrequire} {\textbf{\textsc{Entrées:}}}
%\renewcommand{\algorithmicensure}  {\textbf{\textsc{Sorties:}}}
%\renewcommand{\algorithmicwhile}   {\textbf{tant que}}
%\renewcommand{\algorithmicdo}      {\textbf{faire}}
%\renewcommand{\algorithmicendwhile}{\textbf{fin tant que}}
%\renewcommand{\algorithmicend}     {\textbf{fin}}
%\renewcommand{\algorithmicif}      {\textbf{si}}
%\renewcommand{\algorithmicendif}   {\textbf{fin si}}
%\renewcommand{\algorithmicelse}    {\textbf{sinon}}
%\renewcommand{\algorithmicthen}    {\textbf{alors}}
%\renewcommand{\algorithmicfor}     {\textbf{pour}}
%\renewcommand{\algorithmicforall}  {\textbf{pour tout}}
%\renewcommand{\algorithmicdo}      {\textbf{faire}}
%\renewcommand{\algorithmicendfor}  {\textbf{fin pour}}
%\renewcommand{\algorithmicloop}    {\textbf{boucler}}
%\renewcommand{\algorithmicendloop} {\textbf{fin boucle}}
%\renewcommand{\algorithmicrepeat}  {\textbf{répéter}}
%\renewcommand{\algorithmicuntil}   {\textbf{jusqu'à}}
%\renewcommand{\algorithmiccomment} {\STATE //}
%\newcommand{\BEGIN}{\STATE \fbox{Début}}
%\newcommand{\END}{\STATE \fbox{Fin}}
%\floatname{algorithm}{Algorithme}
%% /francisation des algorithmes

\renewcommand{\qedsymbol}{}

\newcommand{\petit}[1]{
	\medskip \noindent
	\begin{small}
	#1)
	\end{small}
}

\begin{document}
\title{Introduction à Git}
\author{\includegraphics{clubnix}}
\date{\url{https://clubnix.fr}}

\maketitle
%% Laisse page blanche pour verso page de garde %%

\vfill
\pagebreak

%\tableofcontents
\newpage
\strut\thispagestyle{empty}
\vfill
\pagebreak
\tableofcontents
\strut\thispagestyle{empty}
%\setcounter{page}{0}
\newpage
\setcounter{page}{1}

\part{Introduction}

\section{Les gestionnaires de versions}

\paragraph{} Un gestionnaire de version est un outil permettant de maintenir
plusieurs versions, c'est-à-dire qu'il est capable d'enregistrer les différents
états d'un ou de plusieurs fichiers au cours du temps et de naviguer facilement
entres ces états.

\paragraph{} Il en existe plusieurs, les plus connus étant \emph{CVS},
\emph{SVN} et \emph{Git}.  Nous nous intéresserons uniquement à ce dernier qui
est notamment utilisé pour la réalisation de Linux.

\section{Git, un gestionnaire décentralisé}

\paragraph{} Git à la particularité, contrairement à d'autres gestionnaires de
version d'être décentralisé, c'est à dire qu'il n'a pas besoin d'un serveur
pour fonctionner.  Cependant pour des raisons de simplicité il est souvent
utilisé de cette manière.
%TODO ref plus loin

\paragraph{} Chaque personne possède un dépôt privé, et un dépôt public.  Le
dépôt privé sert à faire du développement local, et le public sert à envoyer
des versions du code source afin que d'autres personnes puissent les
télécharger.  En donnant à tout un groupe le même dépôt public, nous obtenons
un système centralisé.

\part{Création de dépôt local}

\section{Clonage}

\paragraph{} La manière la plus fréquente de créer un dépôt
(\emph{repository}) Git est le ``clonage''~: la copie entière d'un projet géré
avec Git. Si l'on prends par exemple le projet de ce cours même qui est stocké
sur \href{https://github.com}{GitHub} à l'adresse
\url{https://github.com/ClubNix/git}, il suffit de taper la commande suivante
pour copier le projet entier dans le dossier courant~:

\begin{lstlisting}
$ git clone https://github.com/ClubNix/git
\end{lstlisting}

\paragraph{} Il aussi possible de manière facultative de spécifier dans quel
dépôt le clonage va se faire~:

\begin{lstlisting}
$ git clone https://github.com/ClubNix/git dossier
\end{lstlisting}

\section{Création d'un dépôt vide}

\paragraph{} Il est bien sûr possible avec Git de créer un dépôt vide. Pour
cela, il suffit de faire~:

\begin{lstlisting}
$ git init
\end{lstlisting}

\paragraph{} Cela va tout simplement ajouter les fichiers propres à Git dans le
répertoire courant et ainsi d'initialiser la gestion de version.

%TODO biblio
%
% http://git-scm.com <- réference
% http://rogerdudler.github.io/git-guide/
% http://stackoverflow.com/questions/315911/git-for-beginners-the-definitive-practical-guide#320140
\end{document}
% vim: spell : spelllang=fr
